\chapter{Apports, limites et perspectives}\label{ch:apports}

\section{Apports}

Le travail apporte les éléments suivants au simulateur. Les engagements matière sont
désormais explicites, et le stock physique est séparé du stock libre. Les approvisionnements en
transit sont représentés comme des réceptions ouvertes, et le crédit du stock passe par un point unique,
ce qui supprime les doubles crédits. MTO et DDMRP se coordonnent par l'Open Supply : une commande
n'ouvre pas d'approvisionnement redondant lorsqu'une réception ouverte suffit à la couvrir. Le mode
MTS historique reste compatible, comme le montre la non-régression. Les matières pilotées par DDMRP
ne sont plus soumises à la politique (s,Q). Enfin, l'ensemble est observable par les journaux, le
tableur et l'export sémantique.

\section{Limites}

Ce travail implémente un noyau DDMRP de gestion matière. Il ne reproduit pas toute la méthodologie
DDMRP. Le positionnement stratégique des stocks n'est pas traité. Les paramètres des buffers sont
expérimentaux, et doivent être calibrés sur des données industrielles, à commencer par celles de ZENER.
Certaines règles de dimensionnement, notamment la composante \emph{orderCycle} et la sécurité
rouge, doivent encore être consolidées bibliographiquement. Les valeurs des promesses client sont
des valeurs par défaut à valider métier. L'allocation de l'Open Supply suit un ordre FIFO, sans
optimisation. Les pics de demande sont détectés et signalés, mais ils ne font pas encore l'objet
d'une politique avancée. Les valeurs des fixtures ne doivent jamais être lues comme des paramètres
optimaux pour ZENER.

Une observation du run C illustre une limite de réactivité. La réception GPL est arrivée
physiquement vers $T \approx 556{,}1$, mais la reprise MTO et la réservation n'ont eu lieu que vers
$T \approx 713$. La cause est la programmation de la vérification finale de l'approvisionnement, qui
attendait une échéance ultérieure. Le mécanisme est correct : la matière est réservée avant Make, et
aucune consommation n'a eu lieu sans engagement. Il s'agit donc d'une possibilité d'amélioration, non d'un
échec du test.

\section{Perspectives}

Plusieurs prolongements sont envisagés, par ordre de maturité. La première est la reprise
événementielle : lors de la réception d'un Open Supply alloué, la commande reprend immédiatement au lieu
d'attendre le prochain contrôle programmé. La deuxième est la calibration des paramètres DDMRP avec les
données ZENER, puis l'analyse de sensibilité des facteurs de lead time et de variabilité. La troisième
est la comparaison expérimentale MTS, MTO et MTO avec DDMRP, sur un même jeu de commandes. Les
prolongements suivants sont plus ouverts : passer à des campagnes multi-commandes à plus grande échelle,
étendre le mécanisme à la stratégie ETO, coupler les décisions aux indicateurs SCOR et VSM, préparer
leur exploitation par l'expert VSM, et tester le mécanisme sur d'autres jeux de données, par exemple
DataCo.

\section{Lecture des références}

La bibliographie distingue deux natures d'éléments. Les références \cite{orlicky} et \cite{ptak} portent
les concepts du MRP et du DDMRP. La référence \cite{apics2017} porte le modèle SCOR. Les formules des
zones et la règle de décision, telles qu'implémentées dans SCONTO-SVU, sont décrites au
\cref{ch:ddmrp} avec leur statut : elles relèvent de l'implémentation expérimentale, et non d'une
référence démontrée.
